\documentclass[10pt]{article}

% ---------- Document Identity (update these only) ----------
\newcommand{\DocStage}{}
\newcommand{\DocTitle}{Your Road to Tier One \textbf{SOF}}
\newcommand{\ProgramName}{182-Week Civilian-to-Selection Readiness Pipeline}
\newcommand{\ProgramSubtitle}{}

% ---------- Page ----------
\usepackage[legalpaper,left=0.5in,right=0.5in,top=0.6in,bottom=0.45in]{geometry}

% ---------- Typography ----------
\usepackage[T1]{fontenc}
\usepackage[utf8]{inputenc}
\usepackage{libertinus}
\usepackage{microtype}
\usepackage{parskip}
\usepackage{ragged2e}
\usepackage{textcomp}
\AtBeginDocument{\color{black}}
\AtBeginDocument{\pagecolor{white}}
\AtBeginDocument{\justifying}

% ---------- Landscape pages ----------
\usepackage{pdflscape}

% ---------- Color ----------
\usepackage[table]{xcolor}
\definecolor{StageBlue}{HTML}{1F4E79}
\definecolor{StageBlueDark}{HTML}{163A5A}
\definecolor{LightGray}{HTML}{F2F4F7}
\definecolor{MidGray}{HTML}{D9DEE5}
\definecolor{RuleGray}{HTML}{C7CED8}
\definecolor{HeaderInk}{HTML}{000000}
\definecolor{Ink}{HTML}{000000}

% ---------- Utility ----------
\usepackage{enumitem}
\sloppy
\setlength{\emergencystretch}{2em}

% ---------- Boxes / UI ----------
\usepackage[most]{tcolorbox}

\tcbset{
  charterTitle/.style={
    enhanced,
    colback=StageBlue!4,
    colframe=StageBlue,
    boxrule=1.25pt,
    arc=3pt,
    left=16pt,right=16pt,top=14pt,bottom=14pt,
    borderline north={3.2pt}{0pt}{StageBlue},
    borderline west={4pt}{0pt}{StageBlue},
    borderline south={1.25pt}{0pt}{StageBlue},
    drop shadow={black!25!white},
  },
  charterRule/.style={
    enhanced,
    colback=LightGray,
    colframe=RuleGray,
    boxrule=0.85pt,
    arc=3pt,
    left=12pt,right=12pt,top=10pt,bottom=10pt,
    borderline west={3pt}{0pt}{StageBlue},
  },
  charterBadge/.style={
    enhanced,
    colback=StageBlue,
    coltext=white,
    colframe=StageBlueDark,
    boxrule=0.6pt,
    arc=2pt,
    left=6pt,right=6pt,top=3pt,bottom=3pt,
  }
}
\newtcbox{\Badge}{nobeforeafter,charterBadge}

% ---------- Lists ----------
\setlist[itemize]{leftmargin=1.5em, itemsep=0.25em, topsep=0.25em}
\setlist[enumerate]{leftmargin=1.8em, itemsep=0.25em, topsep=0.25em}

% ---------- TOC ----------
\usepackage{tocloft}
\setcounter{tocdepth}{3}
\setlength{\cftbeforesecskip}{3pt}
\setlength{\cftbeforesubsecskip}{2pt}
\setlength{\cftbeforesubsubsecskip}{0pt}
\renewcommand{\cftsecfont}{\bfseries}
\renewcommand{\cftsecpagefont}{\bfseries}
\renewcommand{\cftsubsecfont}{\normalfont}
\renewcommand{\cftsubsecpagefont}{\normalfont}
\renewcommand{\cftsubsubsecfont}{\normalfont}
\renewcommand{\cftsubsubsecpagefont}{\normalfont}
\setlength{\cftsecindent}{0pt}
\setlength{\cftsubsecindent}{1.25em}
\setlength{\cftsubsubsecindent}{2.8em}
\setlength{\cftsecnumwidth}{2.1em}
\setlength{\cftsubsecnumwidth}{2.8em}
\setlength{\cftsubsubsecnumwidth}{3.6em}

% Remove the built-in TOC heading so only the custom heading is shown
\makeatletter
\renewcommand{\tableofcontents}{\@starttoc{toc}}
\makeatother

% ---------- Header/Footer: page number only ----------
\usepackage{fancyhdr}
\pagestyle{fancy}
\fancyhf{}
\renewcommand{\headrulewidth}{0pt}
\renewcommand{\footrulewidth}{0pt}
\fancyfoot[C]{\thepage}

% ---------- Section formatting ----------
\usepackage{titlesec}

\titleformat{\section}
  {\Large\bfseries\color{Ink}}
  {\thesection}{0.6em}{}
  [\vspace{0.25em}\textcolor{StageBlue}{\rule{\linewidth}{0.8pt}}]

\titleformat{\subsection}
  {\large\bfseries\color{Ink}}
  {\thesubsection}{0.6em}{}

\titleformat{\subsubsection}
  {\normalsize\bfseries\color{Ink}}
  {\thesubsubsection}{0.6em}{}

\titlespacing*{\section}{0pt}{0.95em}{0.35em}
\titlespacing*{\subsection}{0pt}{0.65em}{0.25em}
\titlespacing*{\subsubsection}{0pt}{0.55em}{0.2em}

% ---------- Table engine ----------
\usepackage{tabularray}
\UseTblrLibrary{booktabs}

% ---------- tabularray: GLOBAL table treatment ----------
\SetTblrInner{
  cells = {fg=black, bg=white, valign=t},
}

% ---------- Header helpers ----------
\newcommand{\Hdr}[1]{\shortstack[c]{\strut #1\strut}}

% ---------- Lines ----------
\newcommand{\TitleLine}{\textcolor{StageBlue}{\rule{0.98\linewidth}{0.95pt}}}
\newcommand{\SubTitleLine}{\textcolor{RuleGray}{\rule{0.98\linewidth}{0.65pt}}}

% ---------- Continued on next page ----------
\newcommand{\ContinuedOnNextPageInline}{%
  \par
  \vfill
  \noindent\textcolor{RuleGray}{\rule{\linewidth}{1pt}}\vspace{-2pt}
  \noindent\hfill\textit{Continued on next page}\par\vspace{-10pt}
  \noindent\textcolor{RuleGray}{\rule{\linewidth}{1pt}}
  \clearpage
}

% ---------- PDF hyperlinks ----------
\usepackage[hidelinks]{hyperref}
\hypersetup{
  colorlinks=false,
  pdfborder={0 0 0},
  linktoc=all,
  pdftitle={\DocTitle},
  pdfauthor={\ProgramName}
}

% =========================================================
\begin{document}

\section*{Stage 12.B — Validity Tags: Definitions and Examples}
\phantomsection
\addcontentsline{toc}{section}{Stage 12.B — Validity Tags: Definitions and Examples}

\subsection*{12.B.1 Validity Tag Definitions}
\phantomsection
\addcontentsline{toc}{subsection}{12.B.1 Validity Tag Definitions}

These definitions mirror Stage 1.F and must not be redefined.

\subsubsection*{Session Validity Tags}
\phantomsection

\begingroup
\scriptsize
\setlength{\tabcolsep}{2pt}
\renewcommand{\arraystretch}{1.12}

\begin{longtblr}{
  width=\linewidth,
  rowhead=1,
  colspec = {
    Q[l,wd=0.90in]
    X[l,3.95]
    X[l,2.55]
    X[l,2.55]
  },
  hlines,
  vlines,
  cells = {hsep=6pt, vsep=6pt, halign=l, valign=t},
  row{odd} = {bg=white},
  row{even} = {bg=LightGray},
  row{1} = {bg=StageBlue!15, fg=black, font=\bfseries, halign=l, valign=t},
}
\Hdr{Tag} &
\Hdr{Definition} &
\Hdr{When Used} &
\Hdr{Gate Impact} \\
\textbf{VALID} &
All required session elements completed and recorded in order: Safety self-check, Warm-up, Mobility, Main work, Cardio, Cooldown, Recovery notes, Post-session notes; required fields complete; Cardio meets minimum continuous standard per governance. &
Default tag when execution and documentation are complete and no required element is missing. &
Gate credit: Yes. Admissible session evidence may be used in gate decision windows. \\
\textbf{MODIFIED} &
All required session elements completed and recorded; a risk-reducing substitution occurred that preserves intent and does not increase risk; documentation includes the substitution and reason code. Cardio still meets minimum continuous standard per governance. &
Used when a safe substitution is executed due to environment, pain, time, equipment, illness posture, or safety constraints, and the session remains complete and documented. &
Gate credit: Yes (conditional). Admissible only within Stage 1.F \textbf{MODIFIED} cap posture; must be reason-coded and review-triggering if repeated. \\
\textbf{INVALID} &
Any required session element missing, or required documentation missing, or an undocumented deviation occurred, or Cardio fails the minimum continuous standard per governance. &
Used when execution is incomplete, documentation is incomplete, or required session structure/controls were not met. &
Gate credit: No. Not admissible; cannot be used to justify \textbf{ADVANCE} and cannot support progression claims. \\
\end{longtblr}

\endgroup

\subsubsection*{Test Validity Tags}
\phantomsection

\begingroup
\scriptsize
\setlength{\tabcolsep}{2pt}
\renewcommand{\arraystretch}{1.12}

\begin{longtblr}{
  width=\linewidth,
  rowhead=1,
  colspec = {
    Q[l,wd=0.90in]
    X[l,3.85]
    X[l,2.75]
    X[l,2.50]
  },
  hlines,
  vlines,
  cells = {hsep=6pt, vsep=6pt, halign=l, valign=t},
  row{odd} = {bg=white},
  row{even} = {bg=LightGray},
  row{1} = {bg=StageBlue!15, fg=black, font=\bfseries, halign=l, valign=t},
}
\Hdr{Tag} &
\Hdr{Definition} &
\Hdr{Common Causes} &
\Hdr{Gate Impact} \\
\textbf{VALID} &
Test executed under the applicable Stage 2 protocol posture with required Stage 2.E logging fields complete (ISO 8601 timestamp with offset, Environment Unit \textbf{ID}, Tool Standard, Warm-up posture where required, Evidence Ref where required, conditions notes); no stop-rule violation; measurement method consistent with protocol. &
Proper execution under controlled conditions; complete logging; stable environment/tooling. &
Gate credit: Yes. Admissible test evidence for tier classification and gates when replication posture is satisfied upstream. \\
\textbf{INVALID} &
Any required Stage 2.E field missing (including Environment Unit \textbf{ID} or Evidence Ref when required), protocol rules violated, warm-up not completed when required, unsafe conditions, stop-rule event, or tool/environment identity not recorded. &
Missing Environment Unit \textbf{ID}; missing Evidence Ref; unlogged tool standard; unsafe weather/surface; pool conditions invalidate; protocol violation; abort event. &
Gate credit: No. Not admissible; treat as not yet tested; retest/reslot per Stage 3 protected windows. \\
\end{longtblr}

\endgroup

\newpage

\subsection*{12.B.2 Reason Code Set}
\phantomsection
\addcontentsline{toc}{subsection}{12.B.2 Reason Code Set}

Reason codes are compact families consistent with Stage 1.F posture. Codes are used as reason-code families for \textbf{SESSION-RC} (sessions) and \textbf{TEST-RC} (tests) without inventing new upstream schemas.

\begingroup
\scriptsize
\setlength{\tabcolsep}{2pt}
\renewcommand{\arraystretch}{1.12}

\begin{longtblr}{
  width=\linewidth,
  rowhead=1,
  colspec = {
    Q[l,wd=1.05in]
    X[l,3.10]
    Q[l,wd=1.30in]
    X[l,3.10]
  },
  hlines,
  vlines,
  cells = {hsep=6pt, vsep=6pt, halign=l, valign=t},
  row{odd} = {bg=white},
  row{even} = {bg=LightGray},
  row{1} = {bg=StageBlue!15, fg=black, font=\bfseries, halign=l, valign=t},
}
\Hdr{Reason Code} &
\Hdr{Meaning} &
\Hdr{Applies To} &
\Hdr{Example} \\
\textbf{ENV} &
Environment constraint or unsafe/invalid conditions requiring deviation or invalidation &
\textbf{SESSION-RC}; \textbf{TEST-RC} &
Ice glaze or lightning forces indoor substitution (session) or reslot (test) \\
PAIN &
Pain or mechanics alteration affecting safe execution &
\textbf{SESSION-RC}; \textbf{TEST-RC} &
Gait-altering pain requires substitution or abort and reslot \\
\textbf{EQUIP} &
Equipment limitation or failure preventing planned execution &
\textbf{SESSION-RC}; \textbf{TEST-RC} &
Timer failure; ruck buckle failure; missing required tool identity \\
\textbf{TIME} &
Time or schedule constraint forcing a safe substitution or truncation &
\textbf{SESSION-RC}; \textbf{TEST-RC} &
Session shortened due to verified schedule constraint; test window missed due to time constraint \\
\textbf{ILL} &
Illness or acute symptoms posture (non-diagnostic) affecting safety/validity &
\textbf{SESSION-RC}; \textbf{TEST-RC} &
Fever/systemic symptoms → terminate attempt; reslot; may trigger \textbf{MEDICAL} \textbf{HOLD} \\
\textbf{SAFETY} &
Safety stop or red-flag posture requiring immediate termination/escalation &
\textbf{SESSION-RC}; \textbf{TEST-RC} &
Chest pressure, near-syncope, severe dizziness → \textbf{STOP} \textbf{NOW}; abort; seek evaluation posture \\
\end{longtblr}

\endgroup

\subsection*{12.B.3 Worked Examples}
\phantomsection
\addcontentsline{toc}{subsection}{12.B.3 Worked Examples}

\begingroup
\scriptsize
\setlength{\tabcolsep}{2pt}
\renewcommand{\arraystretch}{1.12}

\begin{longtblr}{
  width=\linewidth,
  rowhead=1,
  colspec = {
    X[l,2.70]
    Q[l,wd=0.85in]
    Q[l,wd=0.95in]
    X[l,3.10]
    X[l,2.35]
    Q[l,wd=1.35in]
  },
  hlines,
  vlines,
  cells = {hsep=6pt, vsep=6pt, halign=l, valign=t},
  row{odd} = {bg=white},
  row{even} = {bg=LightGray},
  row{1} = {bg=StageBlue!15, fg=black, font=\bfseries, halign=l, valign=t},
}
\Hdr{Scenario} &
\Hdr{Tag} &
\Hdr{Reason Code} &
\Hdr{What Must Be Documented} &
\Hdr{Gate Impact} &
\Hdr{Correct Next Step} \\
Outdoor session moved indoors due to ice glaze; all elements completed and logged &
\textbf{MODIFIED} &
\textbf{ENV} &
Which element changed (Cardio modality/environment); why (unsafe footing); substitution details; confirmation all required elements completed and recorded; \textbf{SESSION-RC} noted &
Gate credit: Yes (conditional). Counts only within \textbf{MODIFIED} cap posture. &
Document substitution \\
Session executed but Cardio stopped at 18 minutes due to dizziness; session terminated &
\textbf{INVALID} &
\textbf{SAFETY} &
Which required element failed (Cardio minimum continuous not met); \textbf{STOP} \textbf{NOW} action taken; symptom description (non-diagnostic); \textbf{SESSION-RC} recorded; escalation notes &
Gate credit: No. Not admissible; triggers conservative gate posture. &
\textbf{MEDICAL} \textbf{HOLD} \\
Session executed with pain escalating to mechanics change; switched to lower-risk modality but completed all elements &
\textbf{MODIFIED} &
PAIN &
Pain description and mechanics impact; substitution performed; confirmation risk reduced; all required elements completed; \textbf{SESSION-RC} recorded &
Gate credit: Yes (conditional). Review required if repeated. &
Document substitution \\
Planned main work replaced due to equipment failure; session completed with safe alternative &
\textbf{MODIFIED} &
\textbf{EQUIP} &
Equipment failure description; safe alternative used; all required elements completed; \textbf{SESSION-RC} recorded &
Gate credit: Yes (conditional). &
Document substitution \\
Session logged late from memory with missing required fields (RPE and pain flags absent) &
\textbf{INVALID} &
\textbf{TIME} &
Late entry verified? If not verifiable: record \textbf{MISSING}; identify missing required fields; tag \textbf{INVALID}; \textbf{SESSION-RC} recorded &
Gate credit: No. Inadmissible due to missing required fields. &
\textbf{HOLD} gate \\
Session completed physically but mobility block was skipped and not substituted &
\textbf{INVALID} &
\textbf{TIME} &
Missing required element (Mobility); no substitution; tag \textbf{INVALID}; \textbf{SESSION-RC} recorded &
Gate credit: No. Inadmissible. &
\textbf{HOLD} gate \\
Illness posture: mild systemic symptoms; session downshifted to minimum viable session with all elements logged &
\textbf{MODIFIED} &
\textbf{ILL} &
Illness posture note (non-diagnostic); downshift rationale; all required elements completed; \textbf{SESSION-RC} recorded &
Gate credit: Yes (conditional). May exceed \textbf{MODIFIED} cap only under medical issue posture; triggers review. &
Document substitution \\
Red-flag symptoms during warm-up; session aborted before main work &
\textbf{INVALID} &
\textbf{SAFETY} &
\textbf{STOP} \textbf{NOW} action; symptom description; aborted exposure; missing required elements documented; \textbf{SESSION-RC} recorded &
Gate credit: No. Inadmissible; \textbf{MEDICAL} \textbf{HOLD} posture considered. &
\textbf{MEDICAL} \textbf{HOLD} \\
\end{longtblr}

\endgroup

\newpage

\subsection*{12.B.4 Test Examples}
\phantomsection
\addcontentsline{toc}{subsection}{12.B.4 Test Examples}

\begingroup
\scriptsize
\setlength{\tabcolsep}{2pt}
\renewcommand{\arraystretch}{1.12}

\begin{longtblr}{
  width=\linewidth,
  rowhead=1,
  colspec = {
    X[l,2.85]
    Q[l,wd=0.85in]
    Q[l,wd=0.95in]
    X[l,3.25]
    X[l,2.25]
    Q[l,wd=1.60in]
  },
  hlines,
  vlines,
  cells = {hsep=6pt, vsep=6pt, halign=l, valign=t},
  row{odd} = {bg=white},
  row{even} = {bg=LightGray},
  row{1} = {bg=StageBlue!15, fg=black, font=\bfseries, halign=l, valign=t},
}
\Hdr{Scenario} &
\Hdr{Tag} &
\Hdr{Reason Code} &
\Hdr{What Must Be Documented} &
\Hdr{Gate Impact} &
\Hdr{Correct Next Step} \\
2-mile run time trial attempted on measured route but Route \textbf{ID} not recorded &
\textbf{INVALID} &
\textbf{ENV} &
ISO 8601 timestamp with offset; Environment Unit \textbf{ID} field is \textbf{MISSING}; Tool Standard recorded; warm-up posture; conditions notes; \textbf{TEST-RC} recorded; Evidence Ref recorded or \textbf{MISSING} &
Gate credit: No. Missing Environment Unit \textbf{ID} makes test inadmissible. &
Reslot to next deload/test week \\
Swim test performed but pool length unit not verified; Pool \textbf{ID} recorded but length unit missing in conditions notes &
\textbf{INVALID} &
\textbf{ENV} &
ISO 8601 timestamp with offset; Environment Unit \textbf{ID} (Pool \textbf{ID}); Tool Standard; warm-up Y/N; conditions notes must include pool length and units; \textbf{TEST-RC} recorded; Evidence Ref recorded &
Gate credit: No. Unit uncertainty invalidates comparability and admissibility. &
Retest \\
\textbf{CAL} test completed but required artifact evidence is missing (Evidence Ref \textbf{MISSING}) &
\textbf{INVALID} &
\textbf{EQUIP} &
ISO 8601 timestamp with offset; Environment Unit \textbf{ID} (Machine/Location note as applicable); Tool Standard; warm-up Y/N; Evidence Ref recorded as \textbf{MISSING}; \textbf{TEST-RC} recorded; conditions notes &
Gate credit: No. Missing Evidence Ref makes test inadmissible when required by protocol posture. &
Retest \\
Ruck time trial aborted due to chest pressure; \textbf{STOP} \textbf{NOW} executed &
\textbf{INVALID} &
\textbf{SAFETY} &
ISO 8601 timestamp with offset; Environment Unit \textbf{ID} (Route \textbf{ID}); Tool Standard; warm-up Y/N; Stop/Abort Y with reason; conditions notes; Evidence Ref; \textbf{TEST-RC} recorded &
Gate credit: No. Safety stop invalidates test; escalation posture controls. &
\textbf{MEDICAL} \textbf{HOLD} \\
Underwater test attempted without governance prerequisites (no lifeguard/spotter confirmation) &
\textbf{INVALID} &
\textbf{SAFETY} &
ISO 8601 timestamp with offset; Pool \textbf{ID}; governance prerequisites recorded as not met; Stop/Abort Y; \textbf{TEST-RC} recorded; Evidence Ref; conditions notes &
Gate credit: No. Governance breach triggers safety handling and reslot; may trigger \textbf{MEDICAL} \textbf{HOLD} posture. &
\textbf{MEDICAL} \textbf{HOLD} \\
\end{longtblr}

\endgroup

Reason code coverage note:

\begin{itemize}
\item \textbf{ENV}, PAIN, \textbf{EQUIP}, \textbf{TIME}, \textbf{ILL}, \textbf{SAFETY} are represented at least once across the combined session and test examples. \textbf{TIME} is primarily a session context code in this set; in tests it is used only when a test window is missed or incomplete due to a verified schedule constraint and is still treated as \textbf{INVALID} for gate use.
\end{itemize}

\newpage

\subsection*{12.B.5 Common Misuse Patterns and Controls}
\phantomsection
\addcontentsline{toc}{subsection}{12.B.5 Common Misuse Patterns and Controls}

\begin{itemize}
\item Silent substitutions → \textbf{MUST} tag \textbf{MODIFIED} and document what changed, why, and the safe substitution used; \textbf{SESSION-RC} required.
\item Reconstructed logs → No inference allowed; record \textbf{UNKNOWN}/\textbf{MISSING}; reconstructed entries remain inadmissible for gates.
\item Invalid tests counted toward gates → Gate credit: No; \textbf{INVALID} tests are treated as not yet tested; retest/reslot only.
\item Tool drift unrecorded → Tool Standard must be logged; tool changes must be documented; unlogged tool drift invalidates comparability.
\item Weather/route changed without Environment Unit \textbf{ID} update → Environment Unit \textbf{ID} plus conditions notes required; otherwise test is \textbf{INVALID} for gate use.
\item Hero day bias → enforce trend and replication posture used upstream; \textbf{VALID} evidence only; do not chase single-day peaks.
\item Mislabeling incomplete sessions as \textbf{MODIFIED} → If required elements are missing or required fields are missing, tag \textbf{INVALID}; \textbf{MODIFIED} requires full element completion and documentation.
\item Treating safety stops as “partial credit” → \textbf{STOP} \textbf{NOW} events invalidate attempts; treat as \textbf{INVALID}; apply escalation posture and reslot discipline.
\end{itemize}

\end{document}